\documentclass[tikz,border=8pt]{standalone}
\usepackage{ctex}
\usetikzlibrary{arrows.meta,positioning,calc}
% Shared visual language for LFS architecture diagrams.
\RequirePackage{../../mymath}

\definecolor{LFSBlue}{HTML}{2563EB}
\definecolor{LFSBlueSoft}{HTML}{DBEAFE}
\definecolor{LFSGreen}{HTML}{15803D}
\definecolor{LFSGreenSoft}{HTML}{DCFCE7}
\definecolor{LFSOrange}{HTML}{C2410C}
\definecolor{LFSOrangeSoft}{HTML}{FFEDD5}
\definecolor{LFSRed}{HTML}{B91C1C}
\definecolor{LFSRedSoft}{HTML}{FEE2E2}
\definecolor{LFSPurple}{HTML}{7E22CE}
\definecolor{LFSPurpleSoft}{HTML}{F3E8FF}
\definecolor{LFSGray}{HTML}{475569}
\definecolor{LFSGraySoft}{HTML}{F1F5F9}

\tikzset{
  lfs picture/.style={
    font=\small,
    node distance=8mm and 12mm,
    >=Latex,
    every path/.style={line cap=round, line join=round}
  },
  block/.style={
    draw=LFSBlue,
    fill=LFSBlueSoft,
    rounded corners=2pt,
    line width=0.8pt,
    align=center,
    inner xsep=7pt,
    inner ysep=5pt,
    minimum height=10mm
  },
  compute/.style={
    block,
    draw=LFSPurple,
    fill=LFSPurpleSoft
  },
  storage/.style={
    block,
    draw=LFSGreen,
    fill=LFSGreenSoft
  },
  interface/.style={
    block,
    draw=LFSOrange,
    fill=LFSOrangeSoft
  },
  risk/.style={
    block,
    draw=LFSRed,
    fill=LFSRedSoft
  },
  neutral/.style={
    block,
    draw=LFSGray,
    fill=LFSGraySoft
  },
  decision/.style={
    diamond,
    draw=LFSOrange,
    fill=LFSOrangeSoft,
    line width=0.8pt,
    align=center,
    inner sep=2pt,
    aspect=2
  },
  group/.style={
    draw=LFSGray,
    rounded corners=3pt,
    dashed,
    line width=0.7pt,
    inner sep=6pt
  },
  arrow/.style={
    -{Latex[length=2.4mm,width=1.6mm]},
    draw=LFSGray,
    line width=0.9pt
  },
  data arrow/.style={
    arrow,
    draw=LFSBlue
  },
  control arrow/.style={
    arrow,
    draw=LFSOrange,
    dashed
  },
  residual/.style={
    arrow,
    draw=LFSGreen,
    rounded corners=4pt
  },
  label text/.style={
    font=\scriptsize,
    text=LFSGray,
    fill=white,
    inner sep=1.5pt
  },
  title/.style={
    font=\bfseries,
    text=LFSGray,
    align=center
  }
}


\begin{document}
\begin{tikzpicture}[
  my picture,
  node distance=7mm and 10mm,
  stage/.style={block,text width=31mm,minimum height=11mm},
  endpoint/.style={neutral,text width=17mm,minimum height=11mm}
]

\node[storage,text width=29mm,minimum height=15mm] (artifact) {压缩权重\\与元数据};

\node[interface,text width=30mm,minimum height=11mm, right=14mm of artifact, yshift=20mm]
  (load-restore) {加载时完整恢复};
\node[storage,stage,right=of load-restore] (resident) {驻留高精度权重};
\node[compute,stage,right=of resident] (load-gemm) {浮点矩阵乘法};

\node[interface,text width=30mm,minimum height=11mm, right=14mm of artifact]
  (forward-restore) {每次前向完整恢复};
\node[storage,stage,right=of forward-restore] (temporary) {高精度临时张量};
\node[compute,stage,right=of temporary] (forward-gemm) {浮点矩阵乘法};

\node[interface,text width=30mm,minimum height=11mm, right=14mm of artifact, yshift=-20mm]
  (fused-decode) {分块解码与乘法融合};
\node[storage,stage,right=of fused-decode] (on-chip) {寄存器或片上\\中间状态};
\node[compute,stage,right=of on-chip] (fused-compute) {融合计算};

\node[endpoint,right=13mm of forward-gemm] (output) {输出};

\draw[data arrow] (artifact.east) -- ++(5mm,0) |- (load-restore.west);
\draw[data arrow] (artifact.east) -- (forward-restore.west);
\draw[data arrow] (artifact.east) -- ++(5mm,0) |- (fused-decode.west);
\draw[data arrow] (load-restore) -- (resident) -- (load-gemm);
\draw[data arrow] (forward-restore) -- (temporary) -- (forward-gemm);
\draw[data arrow] (fused-decode) -- (on-chip) -- (fused-compute);
\draw[data arrow] (load-gemm.east) -- ++(5mm,0) |- (output.north);
\draw[data arrow] (forward-gemm) -- (output);
\draw[data arrow] (fused-compute.east) -- ++(5mm,0) |- (output.south);

\node[label text,above=1mm of resident] {加载后长期驻留};
\node[label text,above=1mm of temporary] {每次计算产生完整中间量};
\node[label text,above=1mm of on-chip] {不写回完整恢复张量};

\end{tikzpicture}
\end{document}
